\chapter*{Preface}
Volume~2 met groups, rings and fields one at a time. This book looks at them together, through the maps between them, and finds that
the most important constructions of algebra are the best possible answers to questions about maps. Its hero is an idea due to
Eilenberg, Mac Lane and Yoneda: an object is known completely by the way other objects map to it.

The first three chapters set up the language. Categories, functors and natural transformations make precise what mathematicians always
meant by ``natural''; Yoneda's lemma and adjoint functors turn that language into theorems. The next chapter multiplies matrices and counts the rank of a bilinear map, which is the number of multiplications it needs. Tensor products of modules, exact sequences and
the snake lemma follow.

The next three chapters are commutative algebra, the algebra of polynomial equations. Hilbert's basis theorem and his
Nullstellensatz turn systems of equations into ideals and ideals into geometry; localisation and integral extensions measure
dimension; and Dedekind domains, where unique factorisation fails for numbers but returns for ideals, give the abstract setting of
algebraic number theory.

Two chapters of homological algebra come next. Derived functors, Tor and Ext measure how far constructions are from preserving exact
sequences, and group cohomology classifies group extensions and explains why the carries of school addition are what they are.

The last four chapters are about representations. Semisimple algebras and induced representations give the general theory and
Burnside's theorem on groups of order $p^aq^b$; Young diagrams give every irreducible representation of the symmetric group, with the
hook length formula; Fourier analysis on the symmetric group proves that about $\frac12n\log n$ random transpositions shuffle a deck
of $n$ cards, and not many fewer; and the Peter--Weyl theorem and Weyl's character formula describe the representations of the
unitary groups, whose random elements have eigenvalues spaced like the zeros of the zeta function.

The book assumes the algebra of Volume~2, the Hilbert spaces and compact operators of Volume~7, and the homology and the integration
formula of Volume~8. Each chapter ends with exercises, each with a hint and a complete solution.

The series asks whether pure order and pure chance can be told apart, and why the primes look random yet obey laws. This floor is the hidden frame: categories, primary decomposition, and the abstract Dedekind domain. The next floor takes up what replaces the bell curve when the load lasts a long time. The class group is an object; a later floor will ask whether it is a random group.
